\section{Validation on a labelled corpus}
\label{sec:nexar}

The competition publishes no labels. That is what makes its metric exploitable
from outside, and it is also what makes the corrected protocol of
Section~\ref{sec:protocol} untestable there: every repair in it is stated in
terms of quantities --- a false-positive rate, an annotated collision time, a
fixed operating point --- that the benchmark does not expose. A protocol that
cannot be run is a proposal, not a result.

This section runs it, on a corpus that is independent of the competition and
fully annotated: the Nexar collision-prediction training split of 1,500 dashcam
clips, 750 of which contain a collision and 750 of which are ordinary driving.
Every positive clip carries two annotations, a collision time $t_{\text{ev}}$
and an alert time $t_{\text{al}}$ marking the earliest moment at which a driver
could still have acted. Three things follow that the competition does not
allow. Average precision and area under the curve become measured quantities
rather than a floor recovered by differencing. Earliness is measured to the
annotated collision rather than to the end of the window. And because half the
corpus contains no collision at all, the cost a warning imposes on ordinary
driving becomes a number rather than an argument.

\subsection{Setup}

The video-blind family of Section~\ref{sec:controls} is regenerated for clips
that differ in length and frame rate: each curve is a closed-form expression in
the frame index, and a generator is given only the clip's frame count and frame
rate, so a curve that wanted a label could not read one. The published metric
is implemented exactly as Section~\ref{sec:metric} states it, with
$t_{ai}$ resolved to the annotated collision in each clip's own timebase.

Against these we place a sighted baseline. It is deliberately modest and no
claim here depends on its accuracy; its only job is to read pixels, so that the
protocol has something sighted to separate from the blind. Five motion
statistics are computed per frame at a fixed \SI{7.5}{\hertz} --- bulk and
localised frame difference, frame difference restricted to the lower-centre of
the image, luminance and its change --- and a logistic model over causal
summaries of them predicts whether the current moment lies inside the annotated
actionable window $[t_{\text{al}}, t_{\text{ev}}]$. The split is 600 clips for
training and 300 held out, fixed before any modelling. Every feature at time
$t$ is a function of samples at or before $t$; the property is asserted by
perturbation rather than assumed, by requiring that altering the tail of a clip
leave every earlier feature bit-identical.

Two properties of the corpus are worth stating before any score, because both
bear on the argument. Its clips carry \textbf{36 distinct frame rates} between
\SI{23.6}{fps} and \SI{31.0}{fps}, so an earliness term counted in frames is
summing quantities of different duration within a single dataset. And its
collisions occur between \SI{3.03}{\second} and \SI{56.80}{\second} into their
clips, an \textbf{18.7-fold spread} in the largest earliness any submission can
collect. The scale objection of Section~\ref{sec:why}(a) does not require two
differently windowed corpora to demonstrate; one is enough.

\subsection{The metric behaves the same way, and the corpus leaks}

Table~\ref{tab:nexar-blind} scores the video-blind family on all 1,500 clips. A
constant risk score of $0.51$ collects \SI{19.09}{\second} of TTA@0.5 and
\SI{19.09}{\second} of STTA@0.5 --- \textbf{99.9\% of the \SI{19.10}{\second}
ceiling} the corpus admits --- while scoring exactly $0.500$ on both
discrimination terms, which is what a constant must score. The result of
Section~\ref{sec:leaderboard} is therefore not a property of one organiser's
scorer. It follows from the composition, and it reappears wherever the
composition is used.

One finding here was not anticipated by the analysis of
Section~\ref{sec:why} and could not have been made on the competition corpus.
Clip duration alone separates collisions from ordinary driving at an AUC of
$0.537$; frame count alone at $0.548$; and \textbf{the frame rate of the video
file alone at $0.533$}. A submission that opens the container header, reads how
fast the file was encoded, and never decodes a frame scores above chance.

This is a property of how the corpus was assembled rather than of the metric,
but it has a direct consequence for anyone evaluating on it: a discrimination
score slightly above chance is not evidence that a method has seen anything.
We therefore add a \emph{metadata-only} control --- a model given the file's
frame rate and duration and nothing else --- to the protocol of
Section~\ref{sec:protocol}, alongside the video-blind one. We did not know to
ask for it before running this study, which is itself the argument for running
controls rather than reasoning about them.

\subsection{What each submission earns on held-out clips}

Table~\ref{tab:nexar-held} scores everything on the same 300 held-out clips.
Under the published metric the ordering is inverted with respect to what the
metric is supposed to reward. The sighted model is the only submission that
discriminates at all, at AUC $0.671$ against $0.500$ for every blind curve, and
it earns \SI{0.045}{\second} of TTA. The constant, at chance discrimination,
earns \SI{18.86}{\second}. \textbf{The composition pays blindness 424 times
better than sight.}

The mechanism is the one Section~\ref{sec:why}(c) identifies: neither timing
term reads the magnitude of a score, only whether it stands above $0.5$. A
model trained to be calibrated on a corpus where roughly one frame in fifty is
inside an actionable window produces probabilities that rarely cross a
threshold set at one half, and is penalised for it. A submission that is simply
pinned above the threshold is not.

\subsection{The protocol, tested and found insufficient}

Section~\ref{sec:protocol} makes a specific and falsifiable prediction about
its own two halves: bounding the earliness term normalises it into the unit
interval but \emph{does not} remove the constant as maximiser, and only
conditioning on a fixed operating point of the discrimination metric does.

That prediction holds exactly. Under bounding alone the constant scores
$1.0000$ --- it still takes the whole of the available anticipation, now
expressed as a ratio. Conditioned at a false-positive rate of 10\% it scores
$0.0000$, because there is no threshold at which it detects anything without
detecting everything. Had the constant failed at the bounding step instead,
either the implementation or the analysis of Section~\ref{sec:why} would have
been wrong.

The protocol does not survive the metadata-only control, and it fails under
both readings of its own averaging rule. Read as written --- normalised
earliness averaged over the positives a submission detects --- the
metadata-only control scores a \textbf{perfect $1.0000$} against the sighted
model's $0.0611$, while detecting 14 of 150 positives to the sighted model's
35. Averaged instead over all positives, charging an undetected one zero, it
scores $0.0933$ against $0.0143$. Either way the submission that reads the file
header outscores the one that reads the road, under the protocol proposed to
prevent exactly this.

The reason is structural and we had not seen it. Conditioning uses a
\emph{clip-level} discrimination metric, which asks whether a submission ranks
collision clips above ordinary ones. It cannot ask whether the alarm, within a
clip, is raised at a useful moment. A curve that is flat across a clip passes
clip-level conditioning on whatever discrimination it has, and then collects
maximal earliness on every clip it detects at all, because a flat curve above
threshold has been above threshold since the first frame.

The two readings above are the second defect. Averaging over detected
positives only --- the natural reading of Section~\ref{sec:protocol}, item 2
--- lets a submission buy a perfect score by detecting almost nothing, which is
precisely what the metadata-only control did. Earliness that is never
delivered, because the alarm was never raised, is not earliness; the statistic
must be averaged over all positives. That repair is necessary and, as the
figures above show, nowhere near sufficient.

\subsection{The repair: measure the alarm, not the ranking}

Both defects have the same root. The protocol conditions on a quantity computed
per clip, while the thing a warning system does --- and the burden it imposes
--- happens per frame. The repair is to measure the alarm itself. Partition
every frame of the corpus into two populations:

\begin{itemize}
\item \textbf{actionable}: frames of a collision clip inside
  $[t_{\text{al}}, t_{\text{ev}}]$, the interval the annotators mark as still
  avoidable;
\item \textbf{non-actionable}: every frame of every ordinary-driving clip, plus
  frames of a collision clip before $t_{\text{al}}$.
\end{itemize}

Frames after the collision belong to neither: nothing is being anticipated
there, and charging them would reward a curve for switching off after an event
it cannot know has occurred. Then report two rates at a common threshold ---
\emph{coverage}, the share of actionable time during which the alarm is raised,
and \emph{duty cycle}, the share of non-actionable time during which it is
raised --- and summarise by coverage at a capped duty cycle.

Fig.~\ref{fig:woc} sweeps both over every threshold. The diagonal is the
no-skill line, and the permanent alarm sits on it, at the top right corner: it
covers every actionable moment and alarms through every other one. That corner
is the point the published composite scores at its maximum.

Table~\ref{tab:nexar-held} gives the summary. At a duty cycle capped at 10\%
the entire video-blind family scores $0.0000$, the metadata-only control
$0.0855$, and the sighted model $0.2833$. The ordering is finally the one the
metric was meant to produce, and it is produced by asking a question the
published composite never asks.

The same partition answers in plain units the question Section~\ref{sec:why}
could only pose. A constant above threshold alarms for \textbf{60 minutes in
every hour} of ordinary driving. The linear ramp alarms for 20.5 minutes per
hour, the metadata-only control for 47.7, and the sighted model for
\SI{0.13}{\second}. None of this is visible to the published composite, whose
timing terms are defined only on clips where a collision occurs, so the burden
falls entirely on the discrimination half --- measured here at $0.500$ for
every submission that pays it.
